\documentclass[11pt,a4paper]{article}
\usepackage{fontspec,amsmath,amssymb,amsthm,microtype}
\setmainfont{lmroman10-regular.otf}[BoldFont=lmroman10-bold.otf,ItalicFont=lmroman10-italic.otf,BoldItalicFont=lmroman10-bolditalic.otf]
\setmonofont{lmmono10-regular.otf}[Scale=MatchLowercase]
\usepackage[margin=25mm]{geometry}
\usepackage[unicode,hidelinks]{hyperref}
\setlength{\emergencystretch}{3em}
\setlength{\parskip}{0.4em}
\setlength{\parindent}{0pt}
\allowdisplaybreaks[2]
\raggedbottom
\title{\texorpdfstring{Passing from a chosen initial coupling to an infimum}{Passing from a chosen initial coupling to an infimum}}
\author{LangevinLean mathematical exposition}
\date{}
\begin{document}\maketitle
\noindent Source module: \texttt{Unadjusted\allowbreak{}Langevin/\allowbreak{}Initial\allowbreak{}Coupling.\allowbreak{}lean}. The accompanying HTML pairs this mathematical exposition with the exact Lean statements and records the scope of the formalization.


The analytic estimate first uses the cost of one initial coupled pair. The infimum calculation here replaces that cost by \(W_2(\mu,\nu)\) when all finite-cost couplings give the same output laws. \texttt{Product\allowbreak{}Reference.\allowbreak{}lean} proves this common-law property on concrete product extensions; \texttt{Canonical\allowbreak{}ULA.\allowbreak{}lean} supplies their stationary references.



The mathematical infimum here is the infimum in \texttt{wasserstein2}. The proof works with arbitrary finite-cost couplings and approximates that infimum when necessary. This gives the true initial transport term without postulating an optimal transport plan.


\section{\texorpdfstring{A uniform affine estimate passes to the infimum}{A uniform affine estimate passes to the infimum}}


Let \(a\ge0\) be finite, \(b\in[0,+\infty]\), and \(\Gamma(\mu,\nu)\ne\varnothing\). If the fixed output measures satisfy

\[W_2(\mu',\nu')\le ac(\gamma)+b\quad\text{for every }\gamma\in\Gamma(\mu,\nu),\]

then \(W_2(\mu',\nu')\le aW_2(\mu,\nu)+b\).



Use the nonempty coupling family and expand transport distances as infima.



Multiplication by finite \(a\ge0\) and addition of \(b\) commute with the nonempty extended-real infimum, including \(a=0\). Thus \(aW_2(\mu,\nu)+b=\inf_\gamma(ac(\gamma)+b)\).



The output distance lies below every member of this last family, hence below its infimum.



{\raggedright\noindent Lean declarations: \texttt{wasserstein2\_\allowbreak{}affine\_\allowbreak{}bound\_\allowbreak{}of\_\allowbreak{}all\_\allowbreak{}couplings}.\par}

\section{\texorpdfstring{A finite infimum supplies a finite-\allowbreak{}cost coupling}{A finite infimum supplies a finite-cost coupling}}


If \(W_2(\mu,\nu)<+\infty\), then some \(\gamma\in\Gamma(\mu,\nu)\) has \(c(\gamma)<+\infty\). Attainment of the infimum is not asserted.



The coupling-cost infimum is strictly below \(+\infty\) by the finiteness hypothesis.



An infimum strictly below a threshold has a member strictly below that threshold. Apply this with threshold \(+\infty\).



{\raggedright\noindent Lean declarations: \texttt{exists\_\allowbreak{}finite\allowbreak{}Cost\_\allowbreak{}coupling}.\par}

\section{\texorpdfstring{The real uniform-\allowbreak{}coupling estimate}{The real uniform-coupling estimate}}


Let \(a,b\ge0\) be finite and assume \(W_2(\mu,\nu)<+\infty\). If \(W_2(\mu',\nu')\le ac(\gamma)+b\) for every initial coupling, then

\[W_2^{\mathbb R}(\mu',\nu')\le aW_2^{\mathbb R}(\mu,\nu)+b.\]

A finite initial distance supplies a coupling. Apply the extended uniform-coupling lemma to obtain the desired inequality in extended reals.



The product \(aW_2(\mu,\nu)\) is finite.



Its sum with finite \(b\) is finite, so the extended output inequality permits monotone real conversion.



Distribute real conversion through the finite sum and product to obtain the real inequality.



{\raggedright\noindent Lean declarations: \texttt{wasserstein2\_\allowbreak{}affine\_\allowbreak{}bound\_\allowbreak{}real\_\allowbreak{}of\_\allowbreak{}all\_\allowbreak{}couplings}.\par}

\section{\texorpdfstring{Finite-\allowbreak{}cost initial couplings suffice}{Finite-cost initial couplings suffice}}


Let \(a,b\ge0\) be finite. Assume both the initial and output transport distances are finite and that every finite-cost \(\gamma\in\Gamma(\mu,\nu)\) satisfies

\[W_2^{\mathbb R}(\mu',\nu')\le a\operatorname{toReal}(c(\gamma))+b.\]

Then \(W_2^{\mathbb R}(\mu',\nu')\le aW_2^{\mathbb R}(\mu,\nu)+b\).



For each finite-cost coupling, embed the real estimate in extended nonnegative reals. Finiteness of the output distance and coupling cost recovers their original extended values, giving \(W_2(\mu',\nu')\le ac(\gamma)+b\).



Apply the all-couplings theorem. For a finite-cost coupling, the needed extended estimate has just been proved.



For an infinite-cost coupling, split into the cases \(a=0\) and \(a\ne0\).



If \(a=0\), choose one finite-cost coupling supplied by the finite initial distance. Its estimate gives \(W_2(\mu',\nu')\le b\), which also supplies the required zero-coefficient estimate for the infinite-cost coupling.



If \(a\ne0\), then \(a>0\), so \(a\cdot(+\infty)+b=+\infty\) and the required inequality is automatic.



{\raggedright\noindent Lean declarations: \texttt{wasserstein2\_\allowbreak{}affine\_\allowbreak{}bound\_\allowbreak{}real\_\allowbreak{}of\_\allowbreak{}finite\_\allowbreak{}couplings}.\par}

\end{document}
